\documentclass[11pt, a4paper]{article}
% ── Shared preamble for all RTOS lab documents ────────────────────────────────
% Usage: % ── Shared preamble for all RTOS lab documents ────────────────────────────────
% Usage: % ── Shared preamble for all RTOS lab documents ────────────────────────────────
% Usage: \input{../preamble}  (from each labNN/ directory)
% ──────────────────────────────────────────────────────────────────────────────

% ── Packages ──────────────────────────────────────────────────────────────────
\usepackage[a4paper, top=2.5cm, bottom=2.5cm, left=2.5cm, right=2.5cm, headheight=28pt]{geometry}
\usepackage[utf8]{inputenc}
\usepackage[T1]{fontenc}
\usepackage{lmodern}
\usepackage{booktabs}
\usepackage{array}
\usepackage{tabularx}
\usepackage{longtable}
\usepackage{multirow}
\usepackage{colortbl}
\usepackage{xcolor}
\usepackage{titlesec}
\usepackage{enumitem}
\usepackage{hyperref}
\usepackage{fancyhdr}
\usepackage{parskip}
\usepackage{listings}
\usepackage[most]{tcolorbox}
\usepackage{fontawesome5}
\usepackage{microtype}
\usepackage{graphicx}
\usepackage{amsmath}

% ── Colors (KMUTNB brand) ──────────────────────────────────────────────────────
\definecolor{kmutnbBlue}{RGB}{0,56,116}
\definecolor{kmutnbGold}{RGB}{186,146,36}
\definecolor{lightGray}{RGB}{245,245,245}
\definecolor{midGray}{RGB}{200,200,200}
\definecolor{codeGray}{RGB}{248,248,248}
\definecolor{codeFrame}{RGB}{220,220,220}
\definecolor{tipteal}{RGB}{0,105,92}
\definecolor{warnAmber}{RGB}{121,85,0}
\definecolor{checkGreen}{RGB}{27,94,32}

% ── Section formatting ─────────────────────────────────────────────────────────
\titleformat{\section}
  {\large\bfseries\color{kmutnbBlue}}
  {\thesection.}{0.5em}{}
  [\color{kmutnbGold}\titlerule]

\titleformat{\subsection}
  {\normalsize\bfseries\color{kmutnbBlue}}
  {\thesubsection.}{0.5em}{}

\titleformat{\subsubsection}
  {\small\bfseries\color{kmutnbBlue}}
  {\thesubsubsection.}{0.5em}{}

\titlespacing*{\section}{0pt}{1.4ex plus .2ex}{0.6ex plus .1ex}
\titlespacing*{\subsection}{0pt}{1.0ex plus .2ex}{0.4ex plus .1ex}

% ── Listings (C and bash) ──────────────────────────────────────────────────────
\lstdefinestyle{cstyle}{
  language        = C,
  basicstyle      = \ttfamily\small,
  keywordstyle    = \bfseries\color{kmutnbBlue},
  commentstyle    = \itshape\color{gray},
  stringstyle     = \color{tipteal},
  numberstyle     = \tiny\color{midGray},
  numbers         = left,
  numbersep       = 6pt,
  stepnumber      = 1,
  backgroundcolor = \color{codeGray},
  frame           = single,
  framerule       = 0.4pt,
  rulecolor       = \color{codeFrame},
  breaklines      = true,
  showstringspaces= false,
  tabsize         = 4,
  xleftmargin     = 1.5em,
  xrightmargin    = 0.5em,
  aboveskip       = 0.8em,
  belowskip       = 0.8em,
}

\lstdefinestyle{bashstyle}{
  language        = bash,
  basicstyle      = \ttfamily\small,
  keywordstyle    = \color{kmutnbBlue},
  commentstyle    = \itshape\color{gray},
  backgroundcolor = \color{black!5},
  frame           = single,
  framerule       = 0.4pt,
  rulecolor       = \color{codeFrame},
  breaklines      = true,
  showstringspaces= false,
  xleftmargin     = 1.5em,
  xrightmargin    = 0.5em,
  aboveskip       = 0.6em,
  belowskip       = 0.6em,
}

\lstdefinestyle{textstyle}{
  basicstyle      = \ttfamily\footnotesize,
  backgroundcolor = \color{black!4},
  frame           = single,
  framerule       = 0.4pt,
  rulecolor       = \color{codeFrame},
  breaklines      = true,
  xleftmargin     = 1.5em,
  xrightmargin    = 0.5em,
  aboveskip       = 0.6em,
  belowskip       = 0.6em,
}

% ── Callout boxes ──────────────────────────────────────────────────────────────
\tcbuselibrary{skins, breakable}

\newtcolorbox{tipbox}[1][]{
  enhanced, breakable,
  colback=tipteal!8, colframe=tipteal, coltitle=white,
  fonttitle=\bfseries\small, title={\faLightbulb\quad Tip#1},
  left=6pt, right=6pt, top=4pt, bottom=4pt,
  boxrule=0.8pt, arc=3pt,
}

\newtcolorbox{warnbox}[1][]{
  enhanced, breakable,
  colback=warnAmber!8, colframe=kmutnbGold, coltitle=white,
  fonttitle=\bfseries\small, title={\faExclamationTriangle\quad Note#1},
  left=6pt, right=6pt, top=4pt, bottom=4pt,
  boxrule=0.8pt, arc=3pt,
}

\newtcolorbox{checkpointbox}[1][]{
  enhanced,
  colback=kmutnbBlue!6, colframe=kmutnbBlue, coltitle=white,
  fonttitle=\bfseries\small, title={\faCheckCircle\quad Checkpoint#1},
  left=6pt, right=6pt, top=4pt, bottom=4pt,
  boxrule=0.8pt, arc=3pt,
}

\newtcolorbox{questionbox}[1][]{
  enhanced, breakable,
  colback=kmutnbGold!10, colframe=kmutnbGold, coltitle=white,
  fonttitle=\bfseries\small, title={\faQuestion\quad Question#1},
  left=6pt, right=6pt, top=4pt, bottom=4pt,
  boxrule=0.8pt, arc=3pt,
}

% ── Checkbox list ──────────────────────────────────────────────────────────────
\newlist{checklist}{itemize}{1}
\setlist[checklist]{label=\faSquare[regular], leftmargin=1.5em, noitemsep}

% ── Misc helpers ───────────────────────────────────────────────────────────────
\newcommand{\file}[1]{\texttt{\small #1}}
\newcommand{\api}[1]{\texttt{\small #1}}

% ── Title block macro ─────────────────────────────────────────────────────────
% Usage: \labtitle{03}{Aperiodic Servers \& Producer--Consumer}{2}{CLO 1 \& CLO 2}
\newcommand{\labtitle}[4]{%
  \begin{center}
    {\color{kmutnbBlue}\rule{\linewidth}{2pt}}\\[6pt]
    {\LARGE\bfseries\color{kmutnbBlue} Laboratory #1}\\[4pt]
    {\large\bfseries #2}\\[2pt]
    {\normalsize Real-Time Operating Systems \enspace\textbullet\enspace
     M.Eng.\ ECE \enspace\textbullet\enspace KMUTNB}\\[8pt]
    {\color{kmutnbGold}\rule{\linewidth}{1pt}}
  \end{center}
  \vspace{0.3cm}
  \begin{tabular}{@{}p{3.8cm} p{10cm}@{}}
    \textbf{Platform}       & NXP FRDM-MCXN236 (ARM Cortex-M33 @ 150\,MHz) \\[2pt]
    \textbf{RTOS}           & FreeRTOS v10.6 \\[2pt]
    \textbf{Toolchain}      & ARM GNU Toolchain (\texttt{arm-none-eabi-gcc}) + Make \\[2pt]
    \textbf{Estimated time} & #3 hours \\[2pt]
    \textbf{CLO mapping}    & #4 \\
  \end{tabular}
  \vspace{0.4cm}
}
  (from each labNN/ directory)
% ──────────────────────────────────────────────────────────────────────────────

% ── Packages ──────────────────────────────────────────────────────────────────
\usepackage[a4paper, top=2.5cm, bottom=2.5cm, left=2.5cm, right=2.5cm, headheight=28pt]{geometry}
\usepackage[utf8]{inputenc}
\usepackage[T1]{fontenc}
\usepackage{lmodern}
\usepackage{booktabs}
\usepackage{array}
\usepackage{tabularx}
\usepackage{longtable}
\usepackage{multirow}
\usepackage{colortbl}
\usepackage{xcolor}
\usepackage{titlesec}
\usepackage{enumitem}
\usepackage{hyperref}
\usepackage{fancyhdr}
\usepackage{parskip}
\usepackage{listings}
\usepackage[most]{tcolorbox}
\usepackage{fontawesome5}
\usepackage{microtype}
\usepackage{graphicx}
\usepackage{amsmath}

% ── Colors (KMUTNB brand) ──────────────────────────────────────────────────────
\definecolor{kmutnbBlue}{RGB}{0,56,116}
\definecolor{kmutnbGold}{RGB}{186,146,36}
\definecolor{lightGray}{RGB}{245,245,245}
\definecolor{midGray}{RGB}{200,200,200}
\definecolor{codeGray}{RGB}{248,248,248}
\definecolor{codeFrame}{RGB}{220,220,220}
\definecolor{tipteal}{RGB}{0,105,92}
\definecolor{warnAmber}{RGB}{121,85,0}
\definecolor{checkGreen}{RGB}{27,94,32}

% ── Section formatting ─────────────────────────────────────────────────────────
\titleformat{\section}
  {\large\bfseries\color{kmutnbBlue}}
  {\thesection.}{0.5em}{}
  [\color{kmutnbGold}\titlerule]

\titleformat{\subsection}
  {\normalsize\bfseries\color{kmutnbBlue}}
  {\thesubsection.}{0.5em}{}

\titleformat{\subsubsection}
  {\small\bfseries\color{kmutnbBlue}}
  {\thesubsubsection.}{0.5em}{}

\titlespacing*{\section}{0pt}{1.4ex plus .2ex}{0.6ex plus .1ex}
\titlespacing*{\subsection}{0pt}{1.0ex plus .2ex}{0.4ex plus .1ex}

% ── Listings (C and bash) ──────────────────────────────────────────────────────
\lstdefinestyle{cstyle}{
  language        = C,
  basicstyle      = \ttfamily\small,
  keywordstyle    = \bfseries\color{kmutnbBlue},
  commentstyle    = \itshape\color{gray},
  stringstyle     = \color{tipteal},
  numberstyle     = \tiny\color{midGray},
  numbers         = left,
  numbersep       = 6pt,
  stepnumber      = 1,
  backgroundcolor = \color{codeGray},
  frame           = single,
  framerule       = 0.4pt,
  rulecolor       = \color{codeFrame},
  breaklines      = true,
  showstringspaces= false,
  tabsize         = 4,
  xleftmargin     = 1.5em,
  xrightmargin    = 0.5em,
  aboveskip       = 0.8em,
  belowskip       = 0.8em,
}

\lstdefinestyle{bashstyle}{
  language        = bash,
  basicstyle      = \ttfamily\small,
  keywordstyle    = \color{kmutnbBlue},
  commentstyle    = \itshape\color{gray},
  backgroundcolor = \color{black!5},
  frame           = single,
  framerule       = 0.4pt,
  rulecolor       = \color{codeFrame},
  breaklines      = true,
  showstringspaces= false,
  xleftmargin     = 1.5em,
  xrightmargin    = 0.5em,
  aboveskip       = 0.6em,
  belowskip       = 0.6em,
}

\lstdefinestyle{textstyle}{
  basicstyle      = \ttfamily\footnotesize,
  backgroundcolor = \color{black!4},
  frame           = single,
  framerule       = 0.4pt,
  rulecolor       = \color{codeFrame},
  breaklines      = true,
  xleftmargin     = 1.5em,
  xrightmargin    = 0.5em,
  aboveskip       = 0.6em,
  belowskip       = 0.6em,
}

% ── Callout boxes ──────────────────────────────────────────────────────────────
\tcbuselibrary{skins, breakable}

\newtcolorbox{tipbox}[1][]{
  enhanced, breakable,
  colback=tipteal!8, colframe=tipteal, coltitle=white,
  fonttitle=\bfseries\small, title={\faLightbulb\quad Tip#1},
  left=6pt, right=6pt, top=4pt, bottom=4pt,
  boxrule=0.8pt, arc=3pt,
}

\newtcolorbox{warnbox}[1][]{
  enhanced, breakable,
  colback=warnAmber!8, colframe=kmutnbGold, coltitle=white,
  fonttitle=\bfseries\small, title={\faExclamationTriangle\quad Note#1},
  left=6pt, right=6pt, top=4pt, bottom=4pt,
  boxrule=0.8pt, arc=3pt,
}

\newtcolorbox{checkpointbox}[1][]{
  enhanced,
  colback=kmutnbBlue!6, colframe=kmutnbBlue, coltitle=white,
  fonttitle=\bfseries\small, title={\faCheckCircle\quad Checkpoint#1},
  left=6pt, right=6pt, top=4pt, bottom=4pt,
  boxrule=0.8pt, arc=3pt,
}

\newtcolorbox{questionbox}[1][]{
  enhanced, breakable,
  colback=kmutnbGold!10, colframe=kmutnbGold, coltitle=white,
  fonttitle=\bfseries\small, title={\faQuestion\quad Question#1},
  left=6pt, right=6pt, top=4pt, bottom=4pt,
  boxrule=0.8pt, arc=3pt,
}

% ── Checkbox list ──────────────────────────────────────────────────────────────
\newlist{checklist}{itemize}{1}
\setlist[checklist]{label=\faSquare[regular], leftmargin=1.5em, noitemsep}

% ── Misc helpers ───────────────────────────────────────────────────────────────
\newcommand{\file}[1]{\texttt{\small #1}}
\newcommand{\api}[1]{\texttt{\small #1}}

% ── Title block macro ─────────────────────────────────────────────────────────
% Usage: \labtitle{03}{Aperiodic Servers \& Producer--Consumer}{2}{CLO 1 \& CLO 2}
\newcommand{\labtitle}[4]{%
  \begin{center}
    {\color{kmutnbBlue}\rule{\linewidth}{2pt}}\\[6pt]
    {\LARGE\bfseries\color{kmutnbBlue} Laboratory #1}\\[4pt]
    {\large\bfseries #2}\\[2pt]
    {\normalsize Real-Time Operating Systems \enspace\textbullet\enspace
     M.Eng.\ ECE \enspace\textbullet\enspace KMUTNB}\\[8pt]
    {\color{kmutnbGold}\rule{\linewidth}{1pt}}
  \end{center}
  \vspace{0.3cm}
  \begin{tabular}{@{}p{3.8cm} p{10cm}@{}}
    \textbf{Platform}       & NXP FRDM-MCXN236 (ARM Cortex-M33 @ 150\,MHz) \\[2pt]
    \textbf{RTOS}           & FreeRTOS v10.6 \\[2pt]
    \textbf{Toolchain}      & ARM GNU Toolchain (\texttt{arm-none-eabi-gcc}) + Make \\[2pt]
    \textbf{Estimated time} & #3 hours \\[2pt]
    \textbf{CLO mapping}    & #4 \\
  \end{tabular}
  \vspace{0.4cm}
}
  (from each labNN/ directory)
% ──────────────────────────────────────────────────────────────────────────────

% ── Packages ──────────────────────────────────────────────────────────────────
\usepackage[a4paper, top=2.5cm, bottom=2.5cm, left=2.5cm, right=2.5cm, headheight=28pt]{geometry}
\usepackage[utf8]{inputenc}
\usepackage[T1]{fontenc}
\usepackage{lmodern}
\usepackage{booktabs}
\usepackage{array}
\usepackage{tabularx}
\usepackage{longtable}
\usepackage{multirow}
\usepackage{colortbl}
\usepackage{xcolor}
\usepackage{titlesec}
\usepackage{enumitem}
\usepackage{hyperref}
\usepackage{fancyhdr}
\usepackage{parskip}
\usepackage{listings}
\usepackage[most]{tcolorbox}
\usepackage{fontawesome5}
\usepackage{microtype}
\usepackage{graphicx}
\usepackage{amsmath}

% ── Colors (KMUTNB brand) ──────────────────────────────────────────────────────
\definecolor{kmutnbBlue}{RGB}{0,56,116}
\definecolor{kmutnbGold}{RGB}{186,146,36}
\definecolor{lightGray}{RGB}{245,245,245}
\definecolor{midGray}{RGB}{200,200,200}
\definecolor{codeGray}{RGB}{248,248,248}
\definecolor{codeFrame}{RGB}{220,220,220}
\definecolor{tipteal}{RGB}{0,105,92}
\definecolor{warnAmber}{RGB}{121,85,0}
\definecolor{checkGreen}{RGB}{27,94,32}

% ── Section formatting ─────────────────────────────────────────────────────────
\titleformat{\section}
  {\large\bfseries\color{kmutnbBlue}}
  {\thesection.}{0.5em}{}
  [\color{kmutnbGold}\titlerule]

\titleformat{\subsection}
  {\normalsize\bfseries\color{kmutnbBlue}}
  {\thesubsection.}{0.5em}{}

\titleformat{\subsubsection}
  {\small\bfseries\color{kmutnbBlue}}
  {\thesubsubsection.}{0.5em}{}

\titlespacing*{\section}{0pt}{1.4ex plus .2ex}{0.6ex plus .1ex}
\titlespacing*{\subsection}{0pt}{1.0ex plus .2ex}{0.4ex plus .1ex}

% ── Listings (C and bash) ──────────────────────────────────────────────────────
\lstdefinestyle{cstyle}{
  language        = C,
  basicstyle      = \ttfamily\small,
  keywordstyle    = \bfseries\color{kmutnbBlue},
  commentstyle    = \itshape\color{gray},
  stringstyle     = \color{tipteal},
  numberstyle     = \tiny\color{midGray},
  numbers         = left,
  numbersep       = 6pt,
  stepnumber      = 1,
  backgroundcolor = \color{codeGray},
  frame           = single,
  framerule       = 0.4pt,
  rulecolor       = \color{codeFrame},
  breaklines      = true,
  showstringspaces= false,
  tabsize         = 4,
  xleftmargin     = 1.5em,
  xrightmargin    = 0.5em,
  aboveskip       = 0.8em,
  belowskip       = 0.8em,
}

\lstdefinestyle{bashstyle}{
  language        = bash,
  basicstyle      = \ttfamily\small,
  keywordstyle    = \color{kmutnbBlue},
  commentstyle    = \itshape\color{gray},
  backgroundcolor = \color{black!5},
  frame           = single,
  framerule       = 0.4pt,
  rulecolor       = \color{codeFrame},
  breaklines      = true,
  showstringspaces= false,
  xleftmargin     = 1.5em,
  xrightmargin    = 0.5em,
  aboveskip       = 0.6em,
  belowskip       = 0.6em,
}

\lstdefinestyle{textstyle}{
  basicstyle      = \ttfamily\footnotesize,
  backgroundcolor = \color{black!4},
  frame           = single,
  framerule       = 0.4pt,
  rulecolor       = \color{codeFrame},
  breaklines      = true,
  xleftmargin     = 1.5em,
  xrightmargin    = 0.5em,
  aboveskip       = 0.6em,
  belowskip       = 0.6em,
}

% ── Callout boxes ──────────────────────────────────────────────────────────────
\tcbuselibrary{skins, breakable}

\newtcolorbox{tipbox}[1][]{
  enhanced, breakable,
  colback=tipteal!8, colframe=tipteal, coltitle=white,
  fonttitle=\bfseries\small, title={\faLightbulb\quad Tip#1},
  left=6pt, right=6pt, top=4pt, bottom=4pt,
  boxrule=0.8pt, arc=3pt,
}

\newtcolorbox{warnbox}[1][]{
  enhanced, breakable,
  colback=warnAmber!8, colframe=kmutnbGold, coltitle=white,
  fonttitle=\bfseries\small, title={\faExclamationTriangle\quad Note#1},
  left=6pt, right=6pt, top=4pt, bottom=4pt,
  boxrule=0.8pt, arc=3pt,
}

\newtcolorbox{checkpointbox}[1][]{
  enhanced,
  colback=kmutnbBlue!6, colframe=kmutnbBlue, coltitle=white,
  fonttitle=\bfseries\small, title={\faCheckCircle\quad Checkpoint#1},
  left=6pt, right=6pt, top=4pt, bottom=4pt,
  boxrule=0.8pt, arc=3pt,
}

\newtcolorbox{questionbox}[1][]{
  enhanced, breakable,
  colback=kmutnbGold!10, colframe=kmutnbGold, coltitle=white,
  fonttitle=\bfseries\small, title={\faQuestion\quad Question#1},
  left=6pt, right=6pt, top=4pt, bottom=4pt,
  boxrule=0.8pt, arc=3pt,
}

% ── Checkbox list ──────────────────────────────────────────────────────────────
\newlist{checklist}{itemize}{1}
\setlist[checklist]{label=\faSquare[regular], leftmargin=1.5em, noitemsep}

% ── Misc helpers ───────────────────────────────────────────────────────────────
\newcommand{\file}[1]{\texttt{\small #1}}
\newcommand{\api}[1]{\texttt{\small #1}}

% ── Title block macro ─────────────────────────────────────────────────────────
% Usage: \labtitle{03}{Aperiodic Servers \& Producer--Consumer}{2}{CLO 1 \& CLO 2}
\newcommand{\labtitle}[4]{%
  \begin{center}
    {\color{kmutnbBlue}\rule{\linewidth}{2pt}}\\[6pt]
    {\LARGE\bfseries\color{kmutnbBlue} Laboratory #1}\\[4pt]
    {\large\bfseries #2}\\[2pt]
    {\normalsize Real-Time Operating Systems \enspace\textbullet\enspace
     M.Eng.\ ECE \enspace\textbullet\enspace KMUTNB}\\[8pt]
    {\color{kmutnbGold}\rule{\linewidth}{1pt}}
  \end{center}
  \vspace{0.3cm}
  \begin{tabular}{@{}p{3.8cm} p{10cm}@{}}
    \textbf{Platform}       & NXP FRDM-MCXN236 (ARM Cortex-M33 @ 150\,MHz) \\[2pt]
    \textbf{RTOS}           & FreeRTOS v10.6 \\[2pt]
    \textbf{Toolchain}      & ARM GNU Toolchain (\texttt{arm-none-eabi-gcc}) + Make \\[2pt]
    \textbf{Estimated time} & #3 hours \\[2pt]
    \textbf{CLO mapping}    & #4 \\
  \end{tabular}
  \vspace{0.4cm}
}


% ── Header / Footer ────────────────────────────────────────────────────────────
\pagestyle{fancy}
\fancyhf{}
\fancyhead[L]{\small\color{kmutnbBlue}\textbf{KMUTNB -- Faculty of Engineering}}
\fancyhead[R]{\small\color{kmutnbBlue}Dept.\ of Electrical and Computer Engineering}
\fancyfoot[L]{\small\color{kmutnbBlue}Real-Time Operating Systems -- M.Eng.\ ECE}
\fancyfoot[C]{\small\thepage}
\fancyfoot[R]{\small\color{kmutnbBlue}Lab 03}
\renewcommand{\headrulewidth}{0.4pt}
\renewcommand{\headrule}{\hbox to\headwidth{%
  \color{kmutnbGold}\leaders\hrule height\headrulewidth\hfill}}
\renewcommand{\footrulewidth}{0.4pt}
\renewcommand{\footrule}{\hbox to\headwidth{%
  \color{midGray}\leaders\hrule height\footrulewidth\hfill}}

\hypersetup{
  colorlinks = true,
  linkcolor  = kmutnbBlue,
  urlcolor   = kmutnbBlue,
  citecolor  = kmutnbBlue,
  pdftitle   = {Lab 03 -- Aperiodic Servers \& Producer--Consumer Pipeline},
  pdfauthor  = {KMUTNB RTOS Course},
}

% ══════════════════════════════════════════════════════════════════════════════
\begin{document}

\labtitle{03}{Aperiodic Servers \& Producer--Consumer Pipeline}{3--4}{CLO 2 \& CLO 3}

% ── Objectives ────────────────────────────────────────────────────────────────
\section{Objectives}

By completing this lab you will be able to:

\begin{enumerate}[noitemsep, leftmargin=*]
  \item Implement a \textbf{deferred-interrupt handler} that moves ISR work into a task context.
  \item Implement a \textbf{Polling Server} pattern using a FreeRTOS software timer and a
        request queue.
  \item Measure \textbf{aperiodic response time} under a periodic background load and compare
        Polling Server vs.\ direct task notification.
  \item Build a \textbf{two-producer, one-consumer} queue pipeline and observe queue
        saturation.
  \item Use an \textbf{event group} to coordinate a multi-condition synchronisation scenario.
\end{enumerate}

% ── Prerequisites ─────────────────────────────────────────────────────────────
\section{Prerequisites}

\begin{checklist}
  \item Lab 02 complete --- you can build, flash, and observe a FreeRTOS application.
  \item ARM GNU Toolchain and \texttt{pyocd} working.
  \item FRDM-MCXN236 board connected via MCU-LINK USB.
  \item SEGGER SystemView installed (needed for Part D).
\end{checklist}

% ── Background ────────────────────────────────────────────────────────────────
\section{Background}

The course Task Model assumes \textbf{periodic} tasks with known execution times. Real systems
also receive \textbf{aperiodic} events (button presses, UART interrupts, sensor alerts) that
must be handled without jeopardising periodic task deadlines.

Two strategies are in scope for this lab:

\renewcommand{\arraystretch}{1.3}
\begin{tabular}{@{} p{3.0cm} p{3.5cm} p{3.5cm} p{3.5cm} @{}}
  \toprule
  \textbf{Strategy} & \textbf{Mechanism} & \textbf{Response time} & \textbf{Periodic impact} \\
  \midrule
  Deferred interrupt &
    ISR gives semaphore $\to$ task unblocks immediately &
    Very low &
    Depends on handler task priority \\
  Polling Server &
    Software timer wakes a service task on a fixed budget and period &
    Up to one server period &
    Bounded --- server is treated as a periodic task \\
  \bottomrule
\end{tabular}
\renewcommand{\arraystretch}{1}

% ── Project Structure ─────────────────────────────────────────────────────────
\section{Project Structure}

\begin{lstlisting}[style=textstyle]
lab03_aperiodic_server/
  Makefile
  src/
    main.c                 <- creates all tasks and kernel objects
    periodic_tasks.c/.h    <- tau1 (100 ms), tau2 (200 ms) background load
    aperiodic_handler.c/.h <- deferred ISR task + Polling Server task
    producer_consumer.c/.h <- two-producer / one-consumer pipeline (Part C)
    FreeRTOSConfig.h
\end{lstlisting}

\begin{warnbox}
  The FreeRTOS kernel and all hardware drivers come from the NXP MCUXpresso SDK ---
  no separate kernel clone is needed. Hardware initialisation (clocks, pin-mux,
  LPUART4) is handled by \texttt{BOARD\_InitHardware()} from the SDK board-support
  package. Use \texttt{PRINTF} from \texttt{fsl\_debug\_console.h} for all console output.
\end{warnbox}

\subsection{Task and Kernel Object Map}

\renewcommand{\arraystretch}{1.25}
\begin{tabular}{@{} p{4.5cm} p{3.5cm} p{5.5cm} @{}}
  \toprule
  \textbf{Object} & \textbf{Type} & \textbf{Notes} \\
  \midrule
  \texttt{vPeriodic1Task}        & Task prio~3, T=100\,ms & Background periodic load \\
  \texttt{vPeriodic2Task}        & Task prio~2, T=200\,ms & Background periodic load \\
  \texttt{vAperiodicHandlerTask} & Task prio~4             & Deferred ISR handler (Part A) \\
  \texttt{vPollingServerTask}    & Task prio~3             & Budget-limited service task (Part B) \\
  \texttt{xAperiodicSem}         & Binary semaphore        & ISR $\to$ handler signal \\
  \texttt{xRequestQueue}         & Queue depth 4           & Polling Server request queue \\
  \texttt{vSensor1Task}          & Task prio~3, T=100\,ms & Producer 1 (Part C) \\
  \texttt{vSensor2Task}          & Task prio~2, T=200\,ms & Producer 2 (Part C) \\
  \texttt{vLoggerTask}           & Task prio~1             & Consumer (Part C) \\
  \texttt{xPipelineQueue}        & Queue depth 10          & Sensor $\to$ Logger (Part C) \\
  \texttt{xPipelineEvents}       & Event group             & \texttt{DATA\_READY} + \texttt{ALARM} bits (Part C) \\
  \bottomrule
\end{tabular}
\renewcommand{\arraystretch}{1}

% ── Part A ─────────────────────────────────────────────────────────────────────
\section{Part A --- Deferred Interrupt Handler}

\subsection{Step 1 --- Wire the Button ISR}

The FRDM-MCXN236 has a user button on GPIO P0\_23. The ISR is already registered in
\file{main.c}; your task is to implement the handler.

In \file{aperiodic\_handler.c}, implement:

\begin{lstlisting}[style=cstyle]
/* ISR -- keep it minimal: signal and return */
void GPIO0_IRQHandler(void)
{
    BaseType_t xHigherPriorityTaskWoken = pdFALSE;
    GPIO_ClearPinsInterruptFlags(GPIO0, 1u << 23);
    xSemaphoreGiveFromISR(xAperiodicSem, &xHigherPriorityTaskWoken);
    portYIELD_FROM_ISR(xHigherPriorityTaskWoken);
}

/* Handler task -- runs at prio 4, higher than all periodic tasks */
void vAperiodicHandlerTask(void *pv)
{
    uint32_t count = 0;
    for (;;) {
        xSemaphoreTake(xAperiodicSem, portMAX_DELAY);
        count++;
        PRINTF("[ASYNC] button press #%u -- handled in task context  t=%u ms\r\n",
               count, (unsigned)xTaskGetTickCount());
    }
}
\end{lstlisting}

\subsection{Step 2 --- Build and Test}

\begin{lstlisting}[style=bashstyle]
cd labs/lab03_aperiodic_server
make
make flash
\end{lstlisting}

Press the user button several times and observe the terminal. Each press should appear as
an \texttt{[ASYNC]} line between the \texttt{[P1]} and \texttt{[P2]} periodic lines.

\begin{checkpointbox}[~A]
  Terminal showing at least 3 button presses interleaved with periodic output,
  with no missed presses.
\end{checkpointbox}

\begin{questionbox}[~A1]
  Press the button rapidly 5 times in under 100\,ms. Do any presses get lost? Why or
  why not?
  (\emph{Hint:} \texttt{xAperiodicSem} is a binary semaphore. What happens to the second
  give if the first has not been taken yet?)
\end{questionbox}

\begin{questionbox}[~A2]
  The handler task runs at priority 4, above all periodic tasks. What is the cost of this
  choice? When would you set the handler priority \emph{below} the periodic tasks?
\end{questionbox}

% ── Part B ─────────────────────────────────────────────────────────────────────
\section{Part B --- Polling Server}

\subsection{Step 1 --- Implement the Server}

A Polling Server is a periodic task with a fixed budget $B_s$. It checks a request queue
and drains up to $B_s$\,ms of work per activation.

In \file{aperiodic\_handler.c}:

\begin{lstlisting}[style=cstyle]
#define POLLING_SERVER_BUDGET_MS  10   /* max service per period */
#define POLLING_SERVER_PERIOD_MS  50   /* Ts */

void vPollingServerTask(void *pv)
{
    TickType_t xLastWake = xTaskGetTickCount();
    for (;;) {
        vTaskDelayUntil(&xLastWake, pdMS_TO_TICKS(POLLING_SERVER_PERIOD_MS));

        TickType_t budget_start = xTaskGetTickCount();
        AperiodicRequest_t req;

        while ((xTaskGetTickCount() - budget_start) <
               pdMS_TO_TICKS(POLLING_SERVER_BUDGET_MS))
        {
            if (xQueueReceive(xRequestQueue, &req, 0) != pdTRUE)
                break;
            PRINTF("[PS] served request type=%d  arrived=%u ms  served=%u ms\r\n",
                   req.type, req.timestamp, (unsigned)xTaskGetTickCount());
        }
    }
}
\end{lstlisting}

\begin{warnbox}
  The Polling Server is parameterised as a periodic task with $C_i = B_s$,
  $T_i = T_s$ for schedulability analysis. Add it to the utilisation calculation ---
  does the total $U$ stay below $\ln 2$?
\end{warnbox}

\subsection{Step 2 --- Submit Requests from an ISR}

Wire the button ISR to enqueue a request instead of giving a semaphore
(switch \file{main.c} to \texttt{\#define LAB\_PART POLLING\_SERVER}):

\begin{lstlisting}[style=cstyle]
void GPIO0_IRQHandler(void)
{
    BaseType_t xHPTW = pdFALSE;
    GPIO_ClearPinsInterruptFlags(GPIO0, 1u << 23);
    AperiodicRequest_t req = { .type = 1,
                               .timestamp = xTaskGetTickCountFromISR() };
    xQueueSendFromISR(xRequestQueue, &req, &xHPTW);
    portYIELD_FROM_ISR(xHPTW);
}
\end{lstlisting}

\subsection{Step 3 --- Observe and Measure}

Press the button and note the time between the press (timestamp logged in the ISR) and
the \texttt{[PS] served} line.

\begin{checkpointbox}[~B]
  Terminal output with at least 5 requests served, annotated with arrival time and
  service time.
\end{checkpointbox}

\begin{questionbox}[~B1]
  What is the worst-case aperiodic response time of the Polling Server in terms of $T_s$?
  With $T_s = 50$\,ms, what is the maximum expected wait? Does your measured data agree?
\end{questionbox}

\begin{questionbox}[~B2]
  You pressed the button 8 times quickly. The queue depth is 4. What happened to the
  other 4 requests? What would change if you used \texttt{xQueueSendFromISR} with
  \texttt{portMAX\_DELAY}?
\end{questionbox}

\begin{questionbox}[~B3]
  Compare the Polling Server with the direct deferred handler from Part~A.
  Fill in this table from your observations:

  \smallskip
  \renewcommand{\arraystretch}{1.3}
  \begin{tabular}{@{} p{4.5cm} p{3.5cm} p{3.5cm} @{}}
    \toprule
    \textbf{Metric} & \textbf{Deferred handler} & \textbf{Polling Server} \\
    \midrule
    Response time (typical) & & \\
    Response time (worst case) & & \\
    Effect on periodic tasks at peak load & & \\
    Analysis complexity & & \\
    \bottomrule
  \end{tabular}
  \renewcommand{\arraystretch}{1}
\end{questionbox}

% ── Part C ─────────────────────────────────────────────────────────────────────
\section{Part C --- Two-Producer, One-Consumer Pipeline}

\subsection{Step 1 --- Build the Pipeline}

Implement in \file{producer\_consumer.c}:

\begin{lstlisting}[style=cstyle]
typedef struct {
    uint8_t  sensor_id;
    int16_t  value;
    uint32_t timestamp;
} SensorReading_t;

void vSensor1Task(void *pv)   /* T=100 ms, prio 3 */
{
    static int16_t raw = 0;
    TickType_t xLastWake = xTaskGetTickCount();
    for (;;) {
        vTaskDelayUntil(&xLastWake, pdMS_TO_TICKS(100));
        raw = (raw + 7) % 1000;
        SensorReading_t r = { .sensor_id = 1, .value = raw,
                              .timestamp = xTaskGetTickCount() };
        if (xQueueSend(xPipelineQueue, &r, 0) != pdTRUE)
            PRINTF("[S1] queue full -- dropped!\r\n");
    }
}

void vSensor2Task(void *pv)   /* T=200 ms, prio 2 */
{
    static int16_t raw = 500;
    TickType_t xLastWake = xTaskGetTickCount();
    for (;;) {
        vTaskDelayUntil(&xLastWake, pdMS_TO_TICKS(200));
        raw = (raw + 13) % 1000;
        SensorReading_t r = { .sensor_id = 2, .value = raw,
                              .timestamp = xTaskGetTickCount() };
        xQueueSend(xPipelineQueue, &r, 0);
    }
}

void vLoggerTask(void *pv)    /* prio 1 -- lowest */
{
    static uint32_t count = 0;
    SensorReading_t r;
    for (;;) {
        xQueueReceive(xPipelineQueue, &r, portMAX_DELAY);
        count++;
        PRINTF("[LOG] #%u  sensor=%d  val=%d  t=%u ms\r\n",
               count, r.sensor_id, r.value, r.timestamp);
        if (r.value > 800)
            xEventGroupSetBits(xPipelineEvents, BIT_ALARM);
        if ((count % 10) == 0)
            xEventGroupSetBits(xPipelineEvents, BIT_DATA_READY);
    }
}
\end{lstlisting}

\subsection{Step 2 --- Add Event Group Coordination}

A fourth task waits for \texttt{BIT\_DATA\_READY} (every 10 items processed) OR
\texttt{BIT\_ALARM} (value $> 800$):

\begin{lstlisting}[style=cstyle]
#define BIT_DATA_READY   (1 << 0)
#define BIT_ALARM        (1 << 1)

void vMonitorTask(void *pv)   /* prio 2 */
{
    for (;;) {
        EventBits_t bits = xEventGroupWaitBits(
            xPipelineEvents,
            BIT_DATA_READY | BIT_ALARM,
            pdTRUE,   /* clear on exit */
            pdFALSE,  /* OR -- either bit suffices */
            portMAX_DELAY);
        if (bits & BIT_DATA_READY)
            PRINTF("[MON] 10-item batch complete  t=%u ms\r\n",
                   (unsigned)xTaskGetTickCount());
        if (bits & BIT_ALARM)
            PRINTF("[MON] ALARM: value exceeded threshold  t=%u ms\r\n",
                   (unsigned)xTaskGetTickCount());
    }
}
\end{lstlisting}

Set \texttt{BIT\_ALARM} inside \texttt{vLoggerTask} when \texttt{r.value > 800}.

\begin{checkpointbox}[~C]
  Terminal showing both \texttt{[MON] 10-item batch} and \texttt{[MON] ALARM} lines
  with correct timing.
\end{checkpointbox}

\begin{questionbox}[~C1]
  Why is the event group AND (\texttt{pdTRUE}) vs.\ OR (\texttt{pdFALSE}) distinction
  important here? What would happen if you used AND for this monitor task?
\end{questionbox}

\begin{questionbox}[~C2]
  Use \texttt{uxQueueMessagesWaiting(xPipelineQueue)} inside \texttt{vLoggerTask} and
  print it every 10 items. Does the queue ever have more than 1 item waiting? Explain
  why or why not at current production vs.\ consumption rates.
\end{questionbox}

\begin{questionbox}[~C3]
  Artificially slow down \texttt{vLoggerTask} with
  \texttt{vTaskDelay(pdMS\_TO\_TICKS(150))} after each print. Watch
  \texttt{[S1] queue full} messages appear. What is the queue high-water mark just
  before the first drop? How does this inform the minimum required queue depth for a
  2-second burst?
\end{questionbox}

% ── Part D ─────────────────────────────────────────────────────────────────────
\section{Part D --- SystemView Trace}

Rebuild without the slowing \texttt{vTaskDelay} from Question~C3. Capture a SystemView
trace of the full system (all tasks running).

\begin{enumerate}[noitemsep, leftmargin=*]
  \item Open SystemView $\to$ Start Recording.
  \item Run for at least 5 seconds.
  \item Stop recording and export the trace.
\end{enumerate}

In the trace identify and annotate:
\begin{itemize}[noitemsep, leftmargin=*]
  \item A context switch from \texttt{vLoggerTask} to \texttt{vSensor1Task} triggered
        by timer expiry.
  \item A \texttt{xQueueSend} call inside \texttt{vSensor1Task} that directly unblocks
        \texttt{vLoggerTask}.
  \item The \texttt{[MON]} task sleeping until the event group fires.
\end{itemize}

\begin{checkpointbox}[~D]
  Annotated screenshot of the SystemView timeline.
\end{checkpointbox}

% ── Deliverables ───────────────────────────────────────────────────────────────
\section{Deliverables}

\renewcommand{\arraystretch}{1.3}
\begin{tabular}{@{} c p{4.5cm} p{7cm} @{}}
  \toprule
  \textbf{\#} & \textbf{Item} & \textbf{Details} \\
  \midrule
  1 & Part A terminal capture   & 3+ button presses interleaved with periodic output \\
  2 & Part B comparison table   & Question~B3 with measured values \\
  3 & Part C terminal capture   & Both event group triggers visible \\
  4 & Part D SystemView screenshot & Annotated timeline \\
  5 & Written answers           & Questions A1--A2, B1--B3, C1--C3 \\
  6 & \file{FreeRTOSConfig.h}   & Utilisation calculation comment including Polling Server \\
  \bottomrule
\end{tabular}
\renewcommand{\arraystretch}{1}

% ── Key Concepts ───────────────────────────────────────────────────────────────
\section{Key Concepts Introduced}

\renewcommand{\arraystretch}{1.3}
\begin{tabular}{@{} p{7cm} p{6cm} @{}}
  \toprule
  \textbf{Concept} & \textbf{Where you saw it} \\
  \midrule
  Deferred interrupt handler pattern         & Part A \\
  Binary semaphore for ISR$\to$task signalling & Part A \\
  Polling Server budget and period           & Part B \\
  Aperiodic response time under periodic load & Parts A \& B \\
  Multi-producer queue pipeline              & Part C \\
  Queue saturation and flow control          & Part C \\
  Event group AND vs.\ OR synchronisation    & Part C \\
  \bottomrule
\end{tabular}
\renewcommand{\arraystretch}{1}

% ── Reference ─────────────────────────────────────────────────────────────────
\section{Reference}

\renewcommand{\arraystretch}{1.3}
\begin{tabular}{@{} p{5.5cm} p{7.5cm} @{}}
  \toprule
  \textbf{Resource} & \textbf{Relevant sections} \\
  \midrule
  Barry, \textit{Mastering the FreeRTOS Kernel} &
    Ch.~6 (semaphores), Ch.~8 (event groups) \\
  Buttazzo, \textit{Hard Real-Time Computing Systems} &
    Ch.~5 (aperiodic task scheduling) \\
  FreeRTOS docs &
    \texttt{xSemaphoreGiveFromISR}, \texttt{xQueueSendFromISR},
    \texttt{xEventGroupSetBits}, \texttt{xEventGroupWaitBits} \\
  \bottomrule
\end{tabular}
\renewcommand{\arraystretch}{1}

\end{document}
